\documentclass[11pt]{article}
\usepackage[margin=1in]{geometry}
\usepackage{amsmath,amssymb}
\usepackage[hidelinks]{hyperref}
\usepackage{enumitem}

\title{Supplementary Information for:\\
Spatial Agreement and Regional Uncertainty in Multi-Source Maps of
\textit{Spartina alterniflora} along the Chinese Coast}

\author{}
\date{}

\begin{document}
\maketitle

This Supplementary Information documents analyses that are superseded by the
results in the main text. They are retained because the underlying result
files remain part of the released experiment record and are referenced by the
main text for traceability. No number in this supplement supports any claim
in the main text.

\section*{S1. Superseded meridian-band regionalisation}
\label{sec:s1-bands}

\subsection*{S1.1 Description}

A preliminary version of the regional analysis split the harmonized
EPSG:3410 domain at $x = 11{,}000{,}000$\,m and $x = 11{,}500{,}000$\,m and
labelled the three slices ``Bohai/Yellow Sea'', ``East China Sea'' and
``South China Sea''. EPSG:3410 is a cylindrical equal-area projection in
which $x$ is linear in longitude (approximately 96.3\,km per degree), so
those splits were meridian bands centred near 114.2\,$^\circ$E and
119.4\,$^\circ$E, not sea regions. The correspondence between the bands and
the labelled seas was inverted: the westernmost band was labelled
``South China Sea'' although it intersects the East China Sea coast, and
the easternmost band was labelled ``Bohai/Yellow Sea'' although it
intersects the South China coast.

All regional statistics in the main text now use IHO-derived marine
polygons (main text, Methods). The band analysis was retained for
traceability under \texttt{results/superseded\_e5\_longitude\_bands/}
in the \textit{Spartina~Earth} repository.

\subsection*{S1.2 Band-shift stress test}

The band-shift experiment (band boundaries shifted by $\pm$250 and
$\pm$500\,km in EPSG:3410) was a stress test of that arbitrary
segmentation, not a geographic sensitivity analysis. Its numbers are
preserved in
\texttt{results/superseded\_e5\_longitude\_bands/e5b\_boundary\_sensitivity.csv}.
They describe sensitivity to meridian-band placement only; they do not
describe Bohai/Yellow Sea, East China Sea or South China Sea as delimited
by IHO marine polygons, and they are not used to support any claim in the
main text. Shifting authoritative marine-polygon coastlines by a fixed
distance, by contrast, is not a meaningful perturbation and was therefore
not performed (main text, Definition sensitivity).

\end{document}
